\section{Affine Macro-Orbit Theory on \texorpdfstring{$\mathbb{Z}_{16}$}{Z16}}
\label{sec:affine-macro-orbits}

\subsection{The Outer Galois Multiplier Group}

Let $\operatorname{Aff}(\mathbb{Z}_n) = \{x \mapsto ax + b \pmod n \mid a \in (\mathbb{Z}_n)^\times, b \in \mathbb{Z}_n\}$ denote the affine group of the cyclic ring $\mathbb{Z}_n$. The dihedral group $D_n = \{x \mapsto \pm x + b \pmod n\}$ is a normal subgroup of $\operatorname{Aff}(\mathbb{Z}_n)$. Consequently, the outer quotient group:
\begin{equation}
\label{eqn:g-macro-def}
G_{\text{macro}} = \operatorname{Aff}(\mathbb{Z}_n) / D_n \cong (\mathbb{Z}_n)^\times / \{\pm 1\}
\end{equation}
acts naturally as a group of permutations on the set of $D_n$-orbits of composition mosaics.

\begin{proposition}[Cyclic Structure of the Macro-Group on $\mathbb{Z}_{16}$]
\label{prop:g-macro-c4}
For $n = 16$, the unit group $(\mathbb{Z}_{16})^\times = \{1, 3, 5, 7, 9, 11, 13, 15\}$ is isomorphic to $C_2 \times C_4$. Factoring out the reflection subgroup $\{\pm 1\} = \{1, 15\}$, the quotient group is cyclic of order $4$:
\begin{equation}
\label{eqn:c4-quotient}
G_{\text{macro}} \cong C_4 = \langle \mu_3 \rangle
\end{equation}
generated by the Galois dilation multiplier:
\begin{equation}
\label{eqn:mu3-def}
\mu_3 \colon \mathbb{Z}_{16} \to \mathbb{Z}_{16}, \quad x \mapsto 3x \pmod{16}
\end{equation}
\end{proposition}

\begin{proof}
The cosets of $(\mathbb{Z}_{16})^\times$ modulo $\{1, 15\}$ are:
\begin{enumerate}
    \item $C_0 = \{1, 15\}$ (identity)
    \item $C_1 = \{3, 13\}$
    \item $C_2 = \{7, 9\}$
    \item $C_3 = \{5, 11\}$
\end{enumerate}
Under multiplication modulo $16$, $3^1 \equiv 3 \in C_1$, $3^2 = 9 \equiv -7 \in C_2$, $3^3 = 27 \equiv 11 \equiv -5 \in C_3$, and $3^4 = 81 \equiv 1 \in C_0$. Thus $C_1$ has order $4$, generating the cyclic quotient $C_4$.
\end{proof}

The unique order-2 subgroup of $G_{\text{macro}}$ is generated by the involution:
\begin{equation}
\label{eqn:mu7-involution}
\mu_7 \colon x \mapsto 7x \equiv -9x \pmod{16}
\end{equation}
which corresponds to the intermediate affine subgroup:
\begin{equation}
\label{eqn:h2-subgroup}
H_2 = D_{16} \rtimes \langle \mu_7 \rangle \le \operatorname{Aff}(\mathbb{Z}_{16})
\end{equation}
of order $64$, whose unit set is $\{1, 7, 9, 15\}$.

We demonstrate such an $x \mapsto 3x \pmod{16}$ inflation with the first block in our example mosaic, $P$, above. This block is the subset $\{0, 2, 7, 9\}$ and is dilated (or inflated) by multiplying each vertex number by $3$, yielding $\{0, 6, 21, 27\}$. The last two elements---having been inflated round the clock, as it were---are not in our ground set except as vertices $5$ and $11$ (via modular arithmetic). So the dilated shape is now the subset $\{0, 5, 6, 11\}$---clearly not the same shape. Dilating again produces $0, 15, 18, 33$---the subset $\{0, 1, 2, 15\}$---and a third dilation yields $0, 3, 6, 45$, i.e.\ $\{0, 3, 6, 13\}$. We know, from the proof above, that a fourth dilation (having come this far without return) \textit{must} return us to the beginning, and $0, 9, 18, 39$---being (modulo $16$) the set $\{0, 2, 7, 9\}$---is indeed that first subset.

Figure~\ref{fig:inflation-in-action} below illustrates the four dilations of the first subset of partition $P$.

\begin{figure}[htbp]
\centering
\input{tkz/T00646-1-dil.tkz}
\caption{\label{fig:inflation-in-action}Successive inflations of subset $\{0, 2, 7, 9\}$}
\end{figure}

The second block of $P$ is $\{1, 4, 10, 11\}$, which maps to numbers $3, 12, 30, 33$---the subset $\{1, 3, 12, 14\}$---which maps again to $3, 9, 36, 42$---the subset $\{3, 4, 9, 10\}$. This maps, for the third time, to $9, 12, 27, 30$---the subset $\{9, 11, 12, 14\}$---which maps for the fourth (and perforce last) time to $27, 33, 36, 42$---the subset $\{1, 4, 10, 11\}$, which is the inevitable second block of $P$ again. The reader should verify that the remaining two blocks, i.e.\ $\{3, 8, 12, 15\}$ and $\{5, 6, 13, 14\}$, perform the same way, and also that at no point in any of the four dilations do the subsets intersect, meaning each stage produces a genuine partition of $\mathbb{Z}_{16}$.

\subsection{Burnside Spectral Decomposition}

Because $G_{\text{macro}} \cong C_4$ acts on the $83{,}488$ dihedral orbits, orbit lengths must divide $|C_4| = 4$. Every affine macro-orbit is therefore one of three types:
\begin{enumerate}
    \item \textbf{Monomers (Size 1)}: Fixed points invariant under all multipliers $\mu_3, \mu_5, \mu_7$.
    \item \textbf{Dimers (Size 2)}: 2-cycles invariant under the involution $\mu_7$, but transposed by $\mu_3 \leftrightarrow \mu_5$.
    \item \textbf{Tetramers (Size 4)}: Full 4-cycles under $\mu_3$:
    \[
    O \xrightarrow{\mu_3} \mu_3(O) \xrightarrow{\mu_3} \mu_7(O) \xrightarrow{\mu_3} \mu_5(O) \xrightarrow{\mu_3} O
    \]
\end{enumerate}

\begin{theorem}[Spectral Decomposition Theorem for $z_4^4$]
\label{thm:spectral-decomposition-z44}
Under the action of $G_{\text{macro}} \cong C_4$, the $83{,}488$ dihedral orbits of shape $z_4^4$ on $\mathbb{Z}_{16}$ decompose into exactly $21{,}030$ Affine Macro-Orbits, comprising:
\begin{itemize}
    \item $c_1 = 14$ monomers (size 1 fixed points)
    \item $c_2 = 295$ dimers (size 2 cycles)
    \item $c_4 = 20{,}721$ tetramers (size 4 cycles)
\end{itemize}
satisfying:
\begin{equation}
\label{eqn:spectral-sum}
c_1 \times 1 + c_2 \times 2 + c_4 \times 4 = 14 + 590 + 82{,}884 = 83{,}488
\end{equation}
and:
\begin{equation}
\label{eqn:total-macro-sum}
c_1 + c_2 + c_4 = 14 + 295 + 20{,}721 = 21{,}030
\end{equation}
\end{theorem}

\begin{proof}
Let $X$ denote the set of $2{,}627{,}625$ raw partitions of $\mathbb{Z}_{16}$ of shape $z_4^4$. The group $\operatorname{Aff}(\mathbb{Z}_{16})$ acts on $X$. The $128$ group elements partition into $9$ cycle classes in the symmetric group $S_{16}$. For each permutation $g \in \operatorname{Aff}(\mathbb{Z}_{16})$, the number of invariant set partitions $|\operatorname{Fix}_X(g)|$ is determined by cycle compatibility between $g$ and the block permutations in $S_4$:
\begin{center}
\begin{tabular}{llrr}
\toprule
Cycle Class of $g \in S_{16}$ & Elements in $\operatorname{Aff}_{16}$ & Fixed Partitions $|\operatorname{Fix}_X(g)|$ \\
\midrule
$(1^{16})$ & $1$ (identity) & $2{,}627{,}625$ \\
$(1^8, 2^4)$ & $2$ & $5{,}425$ \\
$(2^8)$ & $9$ & $4{,}025$ \\
$(1^2, 2^7)$ & $16$ & $945$ \\
$(4^4)$ & $12$ & $89$ \\
$(1^2, 2^3, 4^2)$ & $16$ & $33$ \\
$(1^4, 2^2, 4^2)$ & $8$ & $29$ \\
$(8^2)$ & $32$ & $5$ \\
$(16^1)$ & $32$ & $1$ \\
\bottomrule
\end{tabular}
\end{center}
Applying Burnside's Lemma:
\begin{enumerate}
    \item Dihedral group $D_{16}$ (order $32$, units $\{1, 15\}$):
    \[
    N_{D_{16}} = \frac{1}{32} \sum_{g \in D_{16}} |\operatorname{Fix}_X(g)| = 83{,}488
    \]
    \item Intermediate subgroup $H_2$ (order $64$, units $\{1, 7, 9, 15\}$):
    \[
    N_{H_2} = \frac{1}{64} \sum_{g \in H_2} |\operatorname{Fix}_X(g)| = 42{,}046
    \]
    \item Full affine group $\operatorname{Aff}_{16}$ (order $128$, all $8$ units):
    \[
    N_{\operatorname{Aff}_{16}} = \frac{1}{128} \sum_{g \in \operatorname{Aff}_{16}} |\operatorname{Fix}_X(g)| = 21{,}030
    \]
\end{enumerate}
The orbit counts under $D_{16}$, $H_2$, and $\operatorname{Aff}_{16}$ satisfy the linear system:
\[
\begin{cases}
c_1 + 2c_2 + 4c_4 = 83{,}488 \\
c_1 + 2c_2 + 2c_4 = 42{,}046 \\
c_1 + c_2 + c_4 = 21{,}030
\end{cases}
\]
Subtracting the second equation from the first yields $2c_4 = 83{,}488 - 42{,}046 = 41{,}442$, so $c_4 = 20{,}721$. Subtracting the third equation from the second yields $c_2 + c_4 = 42{,}046 - 21{,}030 = 21{,}016$, giving $c_2 = 21{,}016 - 20{,}721 = 295$. Finally, $c_1 = 21{,}030 - 295 - 20{,}721 = 14$.
\end{proof}
